% ============================================================
% section2_structure.tex - Раздел 2: Инженерное проектирование: Конфигурация, Утилиты и СУБД-архитектура
% ============================================================
% !TeX program = xelatex
% Компиляция: xelatex --shell-escape main.tex

\section{Инженерное проектирование: Конфигурация, Утилиты и СУБД-архитектура}

Прежде чем бросаться писать формулы БПФ, настоящий инженер чертит план. Хаотичный код, где всё свалено в один огромный файл, быстро становится нечитаемым. В этом разделе мы создадим проект в PyCharm, продумаем его модульную структуру, разработаем концепцию перевода физических величин, внедрим СУБД-парадигму хранения данных и напишем наш первый полезный файл — конфигурационный.

\subsection{Создание проекта в PyCharm}

\subsubsection{Новый проект}
Давай создадим рабочее пространство, где будет жить наша программа.

Для начала открой PyCharm, который мы установили в первом разделе. На стартовом экране нажми кнопку \textbf{New Project} (Новый проект).

\smartfigure{fig2.1.png}{Окно создания нового проекта в PyCharm. Важно выбрать правильный базовый интерпретатор и отметить создание виртуального окружения.}{fig:2_1}

В появившемся окне настрой следующие параметры:
\begin{itemize}
    \item \textbf{Location (Расположение):} выбери папку, где будет храниться проект, например: \texttt{C:\textbackslash{}THz\_Project}.
    \item \textbf{Python Interpreter (Интерпретатор):} выбери \texttt{New virtual environment using Virtualenv}. Это автоматически создаст папку \texttt{venv} (наше виртуальное окружение, о котором мы говорили ранее).
    \item Убедись, что выбран базовый интерпретатор Python, который ты установил (например, \texttt{Python 3.12}).
\end{itemize}

Нажми \textbf{Create} (Создать). PyCharm подготовит окружение.

\subsection{Организация файлов}

\subsubsection{Модульная структура}
Представь, что мы собираем сложный физический прибор. Мы не будем паять все детали на одной плате в виде огромного комка проводов. Мы разобьем прибор на функциональные блоки: блок питания, процессор, дисплей.

В программировании это называется \textbf{модульностью}. Мы разделим нашу программу на логические блоки (модули). Каждый модуль (отдельный \texttt{.py} файл) будет решать только свою узкую задачу.

Вот как будет выглядеть архитектура нашего «прибора»:
\begin{itemize}
    \item \texttt{data/} --- папка с исходными текстовыми файлами сигналов (наши «сырые» данные).
    \item \texttt{config.py} --- пульт управления. Здесь мы соберем все настройки и пути к файлам.
    \item \texttt{utils.py} --- модуль утилит. Перевод физических величин и нормализация геометрии.
    \item \texttt{data\_loader.py} --- блок загрузки данных. Читает и структурирует файлы под СУБД-хранилище.
    \item \texttt{spectrum.py} --- блок цифровой обработки (DSP). Удаляет постоянное смещение, накладывает Гауссово окно, вычисляет БПФ без искусственной интерполяции.
    \item \texttt{theoretical.py} --- физико-математический блок. Содержит расчет Blanco-модели и матрицы Джонса.
    \item \texttt{main.py} --- главное интегрированное приложение (GUI) и запуск всего конвейера обработки.
\end{itemize}

Обрати внимание: мы объединили логику визуализации, интерактивной подгонки и главный координационный скрипт в один файл \texttt{main.py}. Это позволяет избежать путаницы между запуском программы и интерфейсными модулями.

Теперь создай в PyCharm все эти файлы (пока пустые) и папку \texttt{data}.

\smartfigure{fig2.2.png}{Дерево файлов проекта в панели Project (слева в PyCharm).}{fig:2_2}

\tip{Чтобы создать файл или папку в PyCharm, кликни правой кнопкой мыши по корневой папке проекта в левой панели (Project) $\rightarrow$ \textbf{New} $\rightarrow$ \textbf{Python File} (или \textbf{Directory}).}

\subsection{Файл конфигурации}

\subsubsection{Зачем нужен config.py}
Любой программный проект начинается с настройки базовых путей к данным. Самая частая ошибка новичка — жестко зашивать пути к файлам на своем жестком диске (например, \texttt{"C:/Users/Student/Desktop/data/..."}) прямо в код загрузчиков. Когда проект переносится на другой компьютер, такой код мгновенно ломается.

Мы поступим профессионально: создадим централизованный файл настроек \texttt{config.py}, который будет определять положение папок относительно самого исполняемого скрипта. Другие модули будут просто импортировать эти настройки.

По мере того, как мы будем развивать физическую модель и добавлять параметры фильтрации или оптимизации, мы будем расширять наш конфигурационный файл, добавляя в него новые константы. Но начнем с самого фундаментального — определения путей.

\subsubsection{Начальная разметка config.py}

Давай напишем код для файла \texttt{config.py}.

\warning{\textbf{Осторожно при копировании из PDF!} Программы для чтения PDF часто «съедают» пробелы в начале строк при копировании. Python строго требует правильных отступов (indentation). Если ты скопируешь код из методички в PyCharm и он подчеркнет его красным, тебе придется вручную восстановить отступы с помощью клавиши \texttt{Tab} или пробелов. Лучший способ научиться — \textbf{перепечатывать код вручную}, так ты быстрее запомнишь синтаксис!}

\begin{minipage}{\linewidth}
	\captionof{listing}{Создаем пульт управления: config.py}
	\begin{minted}{python}
from pathlib import Path

# --- ПУТИ К ДАННЫМ ---
# Используем pathlib для автоматического и платформонезависимого построения путей
BASE_DIR = Path(__file__).resolve().parent
DATA_DIR = BASE_DIR / "data"
	\end{minted}
\end{minipage}

\textbf{Давай разберём, что здесь происходит:}
\begin{itemize}
    \item \texttt{from pathlib import Path}: Модуль \texttt{pathlib} — это современный инструмент Python для работы с путями файлов. Он позволяет удобно строить пути и избегать проблем с обратными слешами в Windows.
    \item \texttt{\_\_file\_\_}: Это специальная переменная Python, которая хранит путь к текущему скрипту (\texttt{config.py}).
    \item \texttt{BASE\_DIR / "data"}: Оператор деления \texttt{/} в \texttt{pathlib} переопределен для удобной «склейки» путей. Это создаст путь к папке \texttt{data} рядом с \texttt{config.py}.
\end{itemize}

\warning{Никогда не используй абсолютные пути вроде \texttt{C:\textbackslash Users\textbackslash ...\textbackslash data}. Если ты перенесешь проект на другой компьютер, путь сломается. Конструкция \texttt{BASE\_DIR / "data"} делает путь \textbf{относительным} — он всегда указывает на папку \texttt{data} внутри текущего проекта.}

\subsection{Разделение уровней абстракции: Модуль утилит utils.py}

Для поддержания чистоты кода и разделения ответственности (Separation of Concerns) мы создаём модуль вспомогательных функций \texttt{utils.py}. Физическое ядро программы должно оперировать строго фундаментальными единицами СИ (метры, секунды, Герцы, линейные амплитуды). Однако пользователю и лаборанту намного привычнее видеть на ползунках микрометры, пикосекунды, децибелы и градусы. 

Модуль \texttt{utils.py} выступает «мостом», изолирующим математику от презентационного слоя. Давай сформулируем первое мини-задание, которое тебе предстоит выполнить.

\begin{minitaskbox}[title={\textbf{🔬 Практическое мини-задание 1: Нормализация диапазона углов ротатора}}]
\mybold{Цель:} Реализовать функцию \texttt{normalize\_angle(angle)}, которая приводит экспериментальные углы с лимба ротатора в диапазоне $[0^\circ, 360^\circ]$ к физически и математически корректному симметричному интервалу $[-180^\circ, 180^\circ]$.

\mybold{Контекст:} В установках ТГц-спектроскопии ротатор поляризатора может иметь шкалу от $0$ до $360$ градусов. При этом углы, превышающие $180^\circ$, фактически эквивалентны отрицательным углам (например, поворот на $280^\circ$ эквивалентен повороту на $280^\circ - 360^\circ = -80^\circ$). Формула Blanco требует именно такого симметричного представления углов относительно оси пропускания.

\mybold{Файл для реализации:} \texttt{utils.py}.
\end{minitaskbox}

Полный листинг модуля \texttt{utils.py} с готовым решением мини-задания приведен в \textbf{Приложении \ref{app:utils}}.

\subsection{СУБД-парадигма: Единое хранилище данных проекта \texttt{data\_store}}

В типовых студенческих скриптах обработка данных часто превращается в хаос: создаются десятки независимых переменных вида \texttt{signals\_raw}, \texttt{background\_signals\_raw}, \texttt{spectra}, \texttt{transmission\_curves}... При этом связь между файлом сигнала образца и соответствующим ему референсным фоном (background) теряется или жестко зашивается через индексы в массивах, что приводит к критическим ошибкам при изменении состава эксперимента.

Чтобы решить эту проблему раз и навсегда, мы применим профессиональный архитектурный паттерн — \textbf{внутрипамяную реляционную базу данных}. Вся информация в нашей программе будет жить в едином централизованном словаре \texttt{data\_store}.

\subsubsection{Схема составных ключей}

Ключом в словаре \texttt{data\_store} является строго типизированный кортеж из четырёх элементов:
\begin{equation}
\text{Key} = (\text{Tag}, \text{Angle}_1, \text{Angle}_2, \text{Repetition})
\end{equation}

Где каждый элемент имеет чёткий физический и логический смысл:
\begin{enumerate}
    \item \textbf{\texttt{Tag} (строка):} Уникальный идентификатор типа данных:
    \begin{itemize}
        \item \texttt{'signal\_raw'} --- сырой временной ТГц-сигнал образца для конкретной реплики (прохода).
        \item \texttt{'bg\_raw'} --- сырой временной референсный сигнал (фон) для конкретного прохода.
        \item \texttt{'spec\_avg'} --- итоговый спектр пропускания по мощности исследуемого образца, полученный путем поштучного вычисления спектра пропускания каждого прохода (отношение мощности сигнала к мощности парного фона) с последующим частотным усреднением по всем повторениям.
        \item \texttt{'spec\_bg\_avg'} --- усредненный спектр мощности опорного сигнала (референса), используемый для контроля динамического диапазона установки.
    \end{itemize}
    \item \textbf{\texttt{Angle1} (число с плавающей точкой):} Нормализованный угол установки первого поляризатора (ротатора) в градусах.
    \item \textbf{\texttt{Angle2} (число с плавающей точкой):} Угол установки второго поляризатора (анализатора).
    \item \textbf{\texttt{Repetition} (целое число или \texttt{None}):} Номер повторного измерения (прохода). Для обработанных, усредненных или спектральных данных, где реплики уже объединены, этот параметр равен \texttt{None}.
\end{enumerate}

\warning{\textbf{Критически важное физическое замечание!}
Студенты часто пытаются сначала усреднить временные ТГц-импульсы (time-domain signals), полученные в нескольких проходах, а затем рассчитать БПФ. \textbf{Делать это категорически запрещено!} 
Из-за температурного расширения элементов, механического люфта линии задержки и микровибраций оптического стола во время эксперимента возникает временной дрейф (так называемый «джиттер» — jitter). Это означает, что ТГц-импульсы в разных проходах могут быть слегка смещены по времени друг относительно друга на фемтосекундные масштабы. 
Если вы усредните их во временной области, небольшие фазовые сдвиги приведут к деструктивной интерференции на высоких частотах. Спектр усредненного сигнала будет содержать искусственное затухание в области 1-2 ТГц. 
Физически корректный подход: рассчитать комплексные спектры для каждого повторения отдельно, взять их модуль (амплитуду или мощность) и усреднять уже \textbf{в частотной области}. Наша структура ключей \texttt{data\_store} прямо форсирует эту логику, храня отдельные сырые сигналы и спектры проходов под соответствующими номерами \texttt{Repetition}.}

\subsubsection{Архитектурные преимущества СУБД-подхода}

\begin{itemize}
    \item \textbf{Абсолютная связность:} При загрузке файлов каждый замер образца автоматически привязывается к своему индивидуальному фону через метаданные угла и номера повторения.
    \item \textbf{Унификация программного интерфейса (API):} Все расчетные функции и графические интерфейсы программы больше не требуют передачи десяти аргументов. Они принимают один единственный объект — \texttt{data\_store}, производят с ним манипуляции и возвращают его же, обогащенного новыми тегами.
    \item \textbf{Элегантная фильтрация (SQL-like):} Сделать выборку нужных данных становится чрезвычайно просто с помощью встроенных в Python генераторов словарей (dict comprehensions).
\end{itemize}

Например, извлечь из базы данных все усредненные спектры образцов для построения графиков можно одной лаконичной строчкой:
\begin{minted}{python}
# Фильтрация по тегу 'spec_avg'
spectra_to_plot = {k: v for k, v in data_store.items() if k[0] == 'spec_avg'}
\end{minted}

Это делает код чрезвычайно читаемым, легко расширяемым и полностью защищенным от случайной перепутанности данных. Чтобы закрепить этот важный архитектурный навык, реши следующее мини-задание.

\begin{minitaskbox}[title={\textbf{🔬 Практическое мини-задание 1b: Расширение структуры СУБД (Кастомный тег)}}]
\mybold{Цель:} Понять внутреннее устройство единой СУБД-структуры проекта, добавив в неё кастомный тег для проведения учебного эксперимента.

\mybold{Постановка задачи:} Студенту предлагается добавить в логику кода кастомный тег (например, \texttt{'signal\_averaged\_time'}). Напиши функцию (или опиши её концептуально), которая для тренировки усредняет ТГц-импульсы прямо во временной области по всем повторениям. Сравни спектры, полученные из такого временного усреднения, с корректными спектрами \texttt{'spec\_avg'}, усредненными в частотной области. Обрати внимание на затухание высоких частот.

\mybold{Файл для реализации:} Студент самостоятельно может реализовать этот эксперимент в виде отдельного учебного скрипта, расширив структуру тегов.
\end{minitaskbox}

\subsection{Полный листинг модуля}

Для обеспечения полноты и удобства полный код конфигурационного модуля \texttt{config.py} вынесен в \textbf{Приложение \ref{app:config}}. На данном этапе в нем расположены только пути к данным, чтобы не загромождать проект неиспользуемыми константами. Мы будем возвращаться к нему и дописывать новые параметры по мере продвижения по книге.

\subsection{Итоги раздела}
Отличный старт! Мы подготовили чертеж нашего будущего приложения:
\begin{enumerate}
    \item Создали изолированный проект со своим виртуальным окружением.
    \item Разбили сложную задачу на логические модули и создали единую базу данных \texttt{data\_store}.
    \item Создали изолированный модуль \texttt{utils.py} для перевода физических величин и нормализации углов.
    \item Поняли физическую опасность временного усреднения сигналов (джиттер) и спроектировали СУБД-хранилище с поддержкой спектрального усреднения.
    \item Объединили логику визуализации и запуска в единый файл \texttt{main.py}, чтобы сделать проект компактным и управляемым.
    \item Создали централизованный конфигурационный файл \texttt{config.py}, содержащий на данном этапе только необходимые файловые пути.
\end{enumerate}

Теперь, когда у нас есть куда складывать данные и откуда брать настройки, самое время научиться эти данные читать. В следующем разделе мы возьмемся за модуль \texttt{data\_loader.py}.
